\section{Axion Cosmology and Dark Matter}
%=============================================================================

%-----------------------------------------------------------------------------
\subsection{Two-Field Structure: Topological vs.\ QCD Axion}
%-----------------------------------------------------------------------------

The TFPT axion sector is split into two distinct fields to avoid an internal
contradiction (a single $\mu$eV axion cannot remain quasi-static over Hubble times):
\begin{equation}
    \boxed{a_{\text{top}} \;\neq\; a_{\text{QCD}}},
    \qquad
    \mathcal{L} \supset \cthree\, a_{\text{top}} F\tilde{F}
    + \frac{a_{\text{QCD}}}{f_a}\,G\tilde{G}.
\end{equation}

\begin{tcolorbox}[colback=orange!5, colframe=orange!60!black, 
    title={\textbf{Two-Field Axion Structure}}, fonttitle=\bfseries]
\begin{itemize}[leftmargin=1.5em, topsep=2pt, itemsep=2pt]
    \item \textbf{$a_{\text{top}}$ (topological field):} Ultralight, quasi-static background
        responsible for birefringence.
        \begin{itemize}[leftmargin=1em, topsep=1pt]
            \item Excursion: $\Delta a_{\text{top}} = n\,\phiz$ (minimal $n=1$)
            \item Time scale: cosmological (Hubble time; no rapid oscillations)
            \item Observable: CMB polarization rotation $\beta = n\,\phiz/(4\pi) = 0.2424^\circ$ for $n=1$
        \end{itemize}
    \item \textbf{$a_{\text{QCD}}$ (QCD axion):} Oscillating field forming dark matter with
        $f_a$ fixed by the PQ block.
        \begin{itemize}[leftmargin=1em, topsep=1pt]
            \item Initial misalignment: $\theta_i \in [0,\pi]$ from $S^1/\mathbb{Z}_2$ field space (see below)
            \item Time scale: $\tau \sim m_a^{-1} \sim 10^{-11}$ s (for $m_a \sim 64\,\mu$eV)
            \item Observable: axion haloscope signal at $\nu \approx 15.6$ GHz
        \end{itemize}
\end{itemize}

\textbf{Minimal mixing:} We adopt the discrete/vanishing mixing limit as the minimal model,
so no new continuous parameters are introduced.
\end{tcolorbox}
\textbf{Why two fields?} A single $\mu$eV axion oscillates at $\sim 10^{10}$ Hz. Over CMB 
propagation time ($\sim 10^{17}$ s), it completes $\sim 10^{27}$ oscillations, averaging 
the birefringence to zero. A distinct ultralight $a_{\text{top}}$ is therefore required.

%-----------------------------------------------------------------------------
\subsection{Dark Matter Relic Density}
%-----------------------------------------------------------------------------

\begin{theorem}{Axion Dark Matter Parameters}{axionDM}
For the QCD axion field $a_{\text{QCD}}$ from the E$_8$ cascade (PQ block at $n = 10$):
\begin{align}
    f_a &= 8.86 \times 10^{10}~\text{GeV} \quad \text{(decay constant)}\\
    m_a &= \frac{m_\pi f_\pi}{f_a}\sqrt{\frac{m_u m_d}{(m_u + m_d)^2}} \approx 64\,\mu\text{eV}\\
    \nu_{\text{haloscope}} &= \frac{m_a}{2\pi\hbar} \approx 15.56~\text{GHz}
\end{align}
\end{theorem}

\textbf{Relic density accounting:}

Including anharmonicity and string-wall production, a standard estimate is
\begin{equation}
    \Omega_a h^2 \simeq 0.12\left(\frac{f_a}{10^{12}~\text{GeV}}\right)^{1.17}\theta_i^2\,F(\theta_i)\,C_{\mathrm{str}},
\end{equation}
where $C_{\mathrm{str}}$ encodes additional production from string-wall networks (order-unity QCD cosmology input, not a TFPT fit parameter).
For $f_a = 8.86 \times 10^{10}$~GeV this gives:
\begin{itemize}[leftmargin=1.5em, topsep=2pt, itemsep=1pt]
    \item For $C_{\mathrm{str}}=1$ (misalignment only), solving yields $\theta_i \approx 2.43$ rad.
    \item \textbf{Suite default (post-inflation PQ):} random $\theta_i$ with $\theta_{\mathrm{rms}}=\pi/\sqrt{3}$ and a \emph{discrete} string/domain-wall enhancement $C_{\mathrm{str}}=7/3$ (minimal cusp/charge set $\{1,2/3,2/3\}$) gives $\Omega_a h^2 \approx 0.123$ (matches $\Omega_{\mathrm{DM}}h^2\simeq 0.12$ at the $\sim 2\%$ level; physics-mode suite gate PASS).
\end{itemize}

\begin{tcolorbox}[colback=yellow!5, colframe=yellow!60!black, 
    title={\textbf{Resolution: Topologically Preferred $\theta_i$}}, fonttitle=\bfseries\small]
\textbf{Option A (topological sector):} On the non-orientable M\"obius fiber, the pseudoscalar
identification enforces $a(x)\mapsto -a(\tau x)$. For constant initial data this implies
$a \equiv -a \;(\mathrm{mod}\;2\pi)$ and hence $a \in \{0,\pi\}$. Equivalently, the field space
is $S^1/\mathbb{Z}_2$, reducing the fundamental domain to $[0,\pi]$ and making $\theta_i$
naturally of order one.

\textbf{Anharmonic enhancement:} Near $\theta_i \to \pi$,
\begin{equation}
    F(\theta_i) \sim \left[\ln\frac{e}{1 - \theta_i/\pi}\right]^{7/6}.
\end{equation}

\textbf{Scenario branches:}
\begin{itemize}[leftmargin=1.5em, topsep=2pt, itemsep=1pt]
    \item \textbf{Pre-inflation PQ (no strings/walls):} $C_{\mathrm{str}}=1$ and a homogeneous $\theta_i$.
    \item \textbf{Post-inflation PQ (strings/walls):} $N_{\mathrm{DW}}=1$ for stability; string-wall networks give $C_{\mathrm{str}}\sim 3$ to $4$.
\end{itemize}

\textbf{Option B (alternative):} Multi-component dark matter:
\begin{itemize}[leftmargin=1.5em, topsep=2pt, itemsep=1pt]
    \item \textbf{Not required in the shipped suite policy:} with the post-inflation RMS branch and discrete $C_{\mathrm{str}}=7/3$, the QCD axion can already saturate the observed DM abundance (torsion-DM remains an optional extension if a different cosmological branch under-produces axions).
\end{itemize}
\end{tcolorbox}

%-----------------------------------------------------------------------------
\subsection{Dark Energy as a Defect-Suppressed Torsion Condensate (candidate)}
\label{sec:dark-energy}
%-----------------------------------------------------------------------------

\begin{tcolorbox}[colback=gray!5, colframe=gray!60,
title={\textbf{Closure-level candidate: discrete $\rho_\Lambda$ with no continuous knobs}}, fonttitle=\bfseries\small]
We test the discrete defect-suppression candidate
\[
\phi_\ast := n\,e^{-\alpha^{-1}(0)/2},
\qquad
\rho_\Lambda \approx (\Mpl \phi_\ast)^4,
\qquad
\Lambda=\rho_\Lambda/\Mpl^2,
\qquad
K_{\mathrm{rms}}=2\sqrt{\Lambda}.
\]
The two-defect sector uses $\delta_2=\tfrac54\delta_{\mathrm{top}}^2$ and the corresponding $\alpha^{-1}(0)$ (Sec.~\ref{sec:sensitivity}).

\textbf{Finite discrete normalization scan:} we scan
\[
n\in\left\{1,\;1/\sqrt{2},\;1/\sqrt{3},\;1/2,\;1/3^{1/4}\right\}
\]
against the $\Lambda$CDM target ledger (from \texttt{k\_calibration.json} in the suite).
The best branch is $n=1/2$, yielding a log$_{10}$ mismatch in $\rho_\Lambda$ of
\[
\Delta\log_{10}\rho_\Lambda \approx 0.086~\text{dex},
\]
within the declared order-of-magnitude falsifiability window (see \codepath{tfpt-suite/2.5/theoryv3/out/dark_energy_exponential_audit}).

\textbf{Why $n=1/2$ is preferred:} the double cover suggests a two-sector split normalization $n=1/2$; this is checked explicitly by \codepath{tfpt-suite/2.5/theoryv3/out/dark_energy_norm_half_origin_audit}.

\textbf{Dynamical completion status:} a spectral-flow-quantized gap-equation version (torsion condensate) exists as a physics-suite module (\texttt{torsion\_condensate}). At the current policy settings it is \emph{close but not closed}: no discrete $n$ hits within $z\le 2$, with best $z\simeq 2.14$ (log$_{10}$ mismatch $\simeq 0.215$ dex). The paper therefore separates (i) the robust normalization audit (PASS candidate) from (ii) the dynamical condensate closure (open interface; FAIL gate in physics mode), to avoid optimistic labeling.
\end{tcolorbox}

%=============================================================================
\section{\texorpdfstring{$\beta_{\text{rad}}$-Centric Identity Summary}{beta_rad-Centric Identity Summary}}
%=============================================================================

The parameter $\beta_{\text{rad}} := \phiz/(4\pi) \approx 0.00423$ provides a compact dictionary for expressing multiple quantities. 
Most equalities in this section are algebraic rewrites of earlier definitions rather than independent new inputs (see Appendix~\ref{app:identities}). 
Inflation observables additionally depend on a dynamical completion (here: the conditional $R^2$ sector in Theorem~\ref{thm:inflation}, assumptions K1--K3).

\begin{table}[H]
\centering
\begin{tabular}{@{}lll@{}}
\toprule
\textbf{Observable} & \textbf{$\beta_{\text{rad}}$ Form} & \textbf{Value}\\
\midrule
Cosmic birefringence & $\beta_{\text{deg}} = \frac{180}{\pi}\,\beta_{\text{rad}}$ & \num[round-mode=places,round-precision=4]{\TFPTbetaDegNum}$^\circ$\\
Neutrino $\theta_{13}$ & $\sin^2\theta_{13} = 4\pi\,\beta_{\text{rad}}\,e^{-\gamma(0)}$ & \num[round-mode=places,round-precision=6]{\TFPTsinSqThetaThirteen}\\
Baryon density (APS seam candidate) & $\Omega_b = (4\pi - 1)\,\beta_{\text{rad}}$ & \num[round-mode=places,round-precision=6]{\TFPTOmegaBPred}\\
\bottomrule
\end{tabular}
\caption{Observables expressed through the single scale $\beta_{\text{rad}} = \phiz/(4\pi)$.}
\label{tab:beta}
\end{table}

\begin{remark}
The $\Omega_b$ identity is no longer a free conjectural coefficient: the APS seam term satisfies $\Delta_\Gamma=2\pi$ (spectral flow equals winding in the minimal nontrivial class).
A natural coefficient candidate is then $K:=2\Delta_\Gamma-1=4\pi-1$, giving $\Omega_b=K\,\beta_{\mathrm{rad}}$.
What remains open is the embedding of the seam operator into the full 4D/6D Dirac operator and a full $\eta$-gluing evaluation on the complete TFPT geometry.
\end{remark}

\TFPTOmegaBTable

%=============================================================================
\section{Summary of Claims and Status}
%=============================================================================

\begingroup
\small
\renewcommand{\arraystretch}{1.15}
\begin{longtable}{@{}p{3cm}p{4cm}p{3cm}p{5.2cm}@{}}
\caption{TFPT status dashboard. ``Exact'': $<0.1\%$; ``Solid'': $<1\%$; ``Minor'': $<3\%$; ``Moderate'': $<5\%$.}\label{tab:status}\\
\toprule
\textbf{Domain} & \textbf{Quantity} & \textbf{Deviation} & \textbf{Status}\\
\midrule
\endfirsthead
\toprule
\textbf{Domain} & \textbf{Quantity} & \textbf{Deviation} & \textbf{Status}\\
\midrule
\endhead
\midrule
\multicolumn{4}{r}{\small Continued on next page.}\\
\midrule
\endfoot
\bottomrule
\endlastfoot
\multicolumn{4}{l}{\textit{Structural (parameter-free)}}\\
\quad CFE & $\alpha^{-1}(0)=\num[round-mode=places,round-precision=12]{\TFPTalphaInvGOverFourNum}$ & $z=\num[round-mode=places,round-precision=3]{\TFPTalphaZGOverFourNum}$ & \textcolor{green!60!black}{\textbf{Derived (double cover + $g=5$ defect partition; no fits)}}\\
\quad UFE & $\beta = \num[round-mode=places,round-precision=4]{\TFPTbetaDegNum}^\circ$ & 0.5--1.7$\sigma$ & \textcolor{green!60!black}{\textbf{Solid}}\\
\quad Cabibbo & $\lambda = \num[round-mode=places,round-precision=6]{0.2243762369}$ & $+0.03\%$ & \textcolor{green!60!black}{\textbf{Solid}}\\
\midrule
\multicolumn{4}{l}{\textit{RG-dressed}}\\
\quad Two-loop & $\alpha_3(1\,\text{PeV})$ & $-1.38\%$ & \textcolor{blue}{\textbf{Good}}\\
\quad Two-loop & $\alpha_3(\mu_{\cthree})$ & $0.00\%$ & \textcolor{green!60!black}{\textbf{Exact}}\\
\quad Unification gate & $\alpha_1,\alpha_2,\alpha_3$ & mismatch $\sim 0.38\%$ & \textcolor{green!60!black}{\textbf{Solid (explicit gate; discrete policy scan)}}\\
\midrule
\multicolumn{4}{l}{\textit{Structural Extensions (Part II)}}\\
\quad PMNS ($\theta_{13}$) & $\sin^2\theta_{13} = \num[round-mode=places,round-precision=5]{0.0231084352}$ & $+5.04\%$ & \textcolor{blue}{\textbf{Testable}}\\
\quad PMNS ($\theta_{12}$) & $\sin^2\theta_{12} = \num[round-mode=places,round-precision=3]{0.318}$ & $+3.56\%$ & \textcolor{blue}{\textbf{Derived}}\\
\quad Inflation ($n_s$) & $n_s \approx 0.9643$ ($N\approx 56$) & within errors & \textcolor{blue}{\textbf{Derived (K1--K3)}}\\
\quad Inflation ($r$) & $r \approx 0.0038$ & future & \textcolor{blue}{\textbf{Derived (K1--K3)}}\\
\quad Bounce transfer function & $T(k)$ large-scale features & future & \textcolor{orange}{\textbf{Conditional (bounce module)}}\\
\quad Axion DM & $m_a \approx 64\,\mu\text{eV}$ & $\theta_i$ model & \textcolor{orange}{\textbf{Structured}}\\
\midrule
\multicolumn{4}{l}{\textit{Suite-closure (physics-mode; assumption-explicit proxy modules)}}\\
\quad Global consistency (suite, physics) & $\chi^2$ scorecard & $p\simeq 0.435$ ($\chi^2\simeq 15.23$, 15 dof) & \textcolor{green!60!black}{\textbf{PASS (dashboard; no full likelihood)}}\\
\quad Flavor (suite, physics) & CKM/PMNS $\chi^2$ & $p_{\mathrm{CKM}}\simeq 0.242$, $p_{\mathrm{PMNS}}\simeq 0.525$ & \textcolor{green!60!black}{\textbf{PASS (imported $\chi^2$ terms)}}\\
\quad Axion DM (relic) & $\Omega_a h^2 \approx 0.123$ & $+2.2\%$ & \textcolor{blue}{\textbf{Derived (post-inflation; $C_{\mathrm{str}}=7/3$)}}\\
\quad Dark energy (norm) & $\rho_\Lambda$ via $e^{-\alpha^{-1}/2}$ & 0.086 dex & \textcolor{green!60!black}{\textbf{Candidate PASS (theoryv3 audit; $n=1/2$)}}\\
\quad Dark energy (condensate) & $\rho_\Lambda$ torsion gap equation & $z\simeq 2.14$ (0.215 dex) & \textcolor{orange}{\textbf{Close, but not closed (FAIL in physics suite)}}\\
\quad BBN & $Y_p, D/H, N_\mathrm{eff}$ & within $2\sigma$ & \textcolor{blue}{\textbf{Consistent (suite)}}\\
\quad Baryogenesis & $\eta_b \approx 5.8\times10^{-10}$ & within 0.5 dex & \textcolor{blue}{\textbf{Proxy (suite)}}\\
\quad Arrow of time & $\log_{10}(S_H)\approx 14.1$ & $>0$ & \textcolor{blue}{\textbf{Proxy (suite)}}\\
\quad GW background & $\Omega_{\rm gw}\approx 3\times10^{-17}$ & $\ll$ bounds & \textcolor{blue}{\textbf{Consistent (suite)}}\\
\quad Precision QED audit & TFPT-scale NP & $\ll$ anomaly scales & \textcolor{blue}{\textbf{Consistent (suite)}}\\
\midrule
\multicolumn{4}{l}{\textit{Requires Further Development}}\\
\quad CKM full matrix & Yukawa textures & RG-dressed baseline + $\chi^2$ diagnostic & \textcolor{blue}{\textbf{Implemented (suite; p-value gate PASS)}}\\
\quad EW scale (constant factory) & $v$ (EW vev) & $320~\mathrm{GeV}$ vs $246~\mathrm{GeV}$ & \textcolor{orange}{\textbf{Pending (candidate; theoryv3 constant factory)}}\\
\quad Lepton ladder (constant factory) & $m_e, m_\mu, m_\tau$ & depends on $v$ candidate & \textcolor{orange}{\textbf{Pending (candidate chain)}}\\
\quad Hadronic scales & $m_p$ placeholder; $m_\pi$ pending & --- & \textcolor{orange}{\textbf{Open (not derived)}}\\
\quad Derived $H_0$ (diagnostic) & from $\Omega_b$ and $\eta_b$ proxy & $z_{\mathrm{Planck}}\simeq -2.99$ & \textcolor{blue}{\textbf{Diagnostic (theoryv3 WARN)}}\\
\quad $\Omega_b$ identity & $(4\pi-1)\beta_{\text{rad}}$ & $z=\num[round-mode=places,round-precision=3]{\TFPTOmegaBZ}$ & \textcolor{blue}{\textbf{Proxy (APS seam candidate; embedding open)}}\\
\end{longtable}
\endgroup

%=============================================================================
\section{What TFPT Does \emph{Not} Assume}
%=============================================================================

The following are \emph{not} required:
\begin{itemize}[leftmargin=2em]
    \item[\ding{55}] Slow-roll inflation (specific model)
    \item[\ding{55}] String landscape or anthropic selection
    \item[\ding{55}] Extra dimensions as free parameters
    \item[\ding{55}] Dark matter entirely from axions
    \item[\ding{55}] New particles beyond minimal axion sector
    \item[\ding{55}] Any fitted or adjusted parameters
\end{itemize}

%=============================================================================
\section{Conclusions}
%=============================================================================

TFPT demonstrates that the fine-structure constant $\alpha$ can be derived 
from first principles using eight minimal axioms rooted in topology, geometry, 
and quantum consistency. The key results are:

\begin{enumerate}
    \item The topological coupling $\cthree = 1/(8\pi)$ from 11D Chern--Simons 
        quantization and the geometric scale $\phiz = 1/(6\pi) + 3/(256\pi^4)$ 
        from M\"obius fiber reduction are the only fundamental inputs.
    
    \item The cubic fixed-point equation (CFE) has a unique positive root. With the derived two-defect term $\delta_2=\frac54\delta_{\mathrm{top}}^2$ (holonomy $g=5$), this yields
        $\alpha^{-1}(0)=\num[round-mode=places,round-precision=12]{\TFPTalphaInvGOverFourNum}$ (one-defect truncation: $\num[round-mode=places,round-precision=10]{\TFPTalphaInvSelfNum}$).
    
    \item The same invariants determine the axion--photon coupling 
        $\gagg = -4\cthree$ and predict cosmic birefringence 
        $\beta = \num[round-mode=places,round-precision=4]{\TFPTbetaDegNum}^\circ$, consistent with observations.
    
    \item An E$_8$ cascade organizes all physical scales from Planck to 
        electroweak, with deviations typically below 5\%.
\end{enumerate}

The theory is falsifiable: any significant revision of $\alpha$, $\beta$, 
or the two-loop RG fingerprints would invalidate the framework.

\paragraph{Outlook (open interfaces).}
Two longer-range directions are sketched in internal notes and remain explicitly open:
an asymptotic-safety UV completion for the torsionful gravity sector and a black-hole
vortex/remnant picture with information-paradox implications. Both are framed as outlook
items rather than closed predictions at this stage \cite{fiveProblemsNote2025,blackHolesBh1Note2025}.

%=============================================================================
\appendix
\section{Numerical Constants (High Precision)}
%=============================================================================

\begin{table}[H]
\centering
\small
\begin{tabular}{@{}ll@{}}
\toprule
\textbf{Quantity} & \textbf{Value}\\
\midrule
$\cthree = 1/(8\pi)$ & \num[round-mode=places,round-precision=12]{\TFPTcThreeNum}\\
$\phitree = 1/(6\pi)$ & \num[round-mode=places,round-precision=12]{\TFPTvarphiTreeNum}\\
$\deltatop = 3/(256\pi^4)$ & \num[round-mode=places,round-precision=12]{\TFPTdeltaTopNum}\\
$\phiz = \phitree + \deltatop$ & \num[round-mode=places,round-precision=12]{\TFPTvarphiZeroNum}\\
$\phiz(\alpha)$ (self-consistent) & \num[round-mode=places,round-precision=12]{0.053170209119023745}\\
$\beta_{\text{rad}} = \phiz/(4\pi)$ & \num[round-mode=places,round-precision=12]{\TFPTbetaRadNum}\\
$\beta_{\text{deg}} = (180/\pi)\,\beta_{\text{rad}}$ & \num[round-mode=places,round-precision=6]{\TFPTbetaDegNum}$^\circ$\\
$g_{a\gamma\gamma}$ & \num[round-mode=places,round-precision=12]{\TFPTgAggNum}\\
$\gamma(0)=g/(g+1)$ (E$_8$) & $\tfrac56$\\
$\lambda$ (E$_8$ exponent) & \num[round-mode=places,round-precision=12]{0.587233190787837868}\\
$M/\Mpl$ ($R^2$ scale) & \num[round-mode=places,round-precision=12]{\TFPTMOverMplNum}\\
$A = 2\cthree^3$ & \num[round-mode=figures,round-precision=8]{1.2598255637968552e-4}\\
$\mathcal{K} = (\bone/2\pi)\ln(1/\phiz)$ & \num[round-mode=places,round-precision=10]{1.91468479457699546}\\
$\alpha^{-1}_{\text{baseline}}$ & \num[round-mode=places,round-precision=9]{137.036501464885822}\\
$\alpha^{-1}_{\text{TFPT}}$ (one-defect) & \num[round-mode=places,round-precision=12]{\TFPTalphaInvSelfNum}\\
$\alpha^{-1}_{\text{TFPT}}$ (two-defect, $g=5$) & \num[round-mode=places,round-precision=12]{\TFPTalphaInvGOverFourNum}\\
$\alpha^{-1}_{\text{TFPT}}$ (two-defect, suite diagnostic) & \num[round-mode=places,round-precision=12]{\TFPTalphaInvTwoDefNum}\\
$\alpha^{-1}_{\text{CODATA 2022}}$ & \num[round-mode=places,round-precision=9]{\TFPTalphaInvCODATANum}\\
Accuracy (ppm; one-defect) & \num[round-mode=places,round-precision=6]{-0.03704}\\
Accuracy (ppm; two-defect, $g=5$) & \num[round-mode=places,round-precision=6]{\TFPTalphaPpmGOverFourNum}\\
Accuracy (ppm; two-defect, suite diagnostic) & \num[round-mode=places,round-precision=6]{\TFPTalphaPpmTwoDefNum}\\
$\xi = \cthree/\phiz$ & \num[round-mode=places,round-precision=9]{0.748303083562547799}\\
\bottomrule
\end{tabular}
\caption{High-precision numerical values of TFPT constants.}
\label{tab:constants}
\end{table}

%=============================================================================
